\chapter{引言}

\section{宇宙学分析的方法与手段}
现代宇宙学分析利用多种天文探针来测量和约束宇宙学参数。星系巡天项目（如DESI）通过精确测量星系光谱，构建星系在三维空间中的分布图，并利用星系的成团性特征来约束宇宙学参数；引力透镜效应则借助宇宙微波背景辐射的畸变（CMB透镜）以及星系在引力场中形态的畸变（弱引力透镜），反推宇宙中的物质分布，从而实现对宇宙学参数的限制；超新星作为标准烛光，其光变曲线的测量同样可用于约束宇宙学参数。

上述不同探针在约束宇宙学参数时，普遍采用贝叶斯分析方法。其基本思路为：将观测到的示踪天体数据压缩为某种数据向量；如果对于给定的宇宙学模型，能够计算出在该模型下理论预期的数据向量及其误差分布（即协方差矩阵），那么贝叶斯分析便可通过采样宇宙学参数，得到其在参数空间中的后验概率分布。

\subsection{协方差矩阵与 N 体模拟}

协方差矩阵是宇宙学模型的重要组成部分，包含统计误差与系统误差两个方面。在某些情况下，协方差矩阵可以通过解析方法计算，但在更多实际问题中，解析表达式难以获得，只能依赖大量模拟进行估计。在宇宙学背景下，对于星系巡天、弱引力透镜等探针，观测到的示踪天体红移较低，引力非线性效应显著，因此必须通过数值模拟来计算协方差矩阵。

宇宙学模拟的主流方法是N体模拟。在冷暗物质模型假设下，根据给定的宇宙学参数生成初始时刻暗物质粒子的分布，然后通过引力相互作用演化这些粒子，最终得到目标红移处的暗物质密度分布和速度场。针对不同探针，还需进一步处理以生成相应示踪天体的模拟分布。

在相同宇宙学参数下，尽管不同模拟实例遵循相同的初始概率分布，但粒子的具体初始位置存在差异，这会导致最终得到的暗物质密度分布和示踪天体分布有所不同，从而产生数据向量的统计误差，这也是协方差矩阵统计误差的主要来源。协方差矩阵中的系统误差部分，有时也可通过将系统误差纳入模拟的方式来估计，但这不在本文的讨论范围之内。

\section{宇宙学快速模拟方法}

为达到宇宙学测量所需的分辨率，N体模拟通常需要Gpc尺度的盒子边长以及数亿甚至数十亿的粒子。如此大规模的粒子在相互引力作用下的演化，将消耗巨大的计算资源，计算时间往往长达数天甚至数月。然而，为了获得较为准确的协方差矩阵估计，通常需要上千次这样的模拟，因此发展快速模拟方法显得尤为重要。

\subsection{宇宙学快速模拟方法的需求}

随着宇宙学观测设备的不断迭代和新一代巡天项目的发展，快速模拟方法也需要同步升级以满足更高的要求。一方面，新一代巡天项目需要更大体积的宇宙学模拟，这意味着需要演化更多的暗物质粒子；另一方面，不同探针对宇宙学参数的响应各不相同，多探针联合分析能够显著提升对宇宙学模型的约束能力。为了正确实现多探针联合分析，必须使用同一套模拟来估计不同探针之间的相关性，即不同探针的协方差矩阵应来源于相同的宇宙学模拟。因此，新的快速模拟方法需要兼顾多种探针的需求。

综上所述，新的快速模拟方法应满足以下几个要求：
\begin{itemize}
  \item 生成速度快，能够在短时间内产生大量用于误差估计的宇宙学模拟；
  \item 易于扩展到大规模模拟，不会因体积增大或粒子数增多而丧失速度优势；
  \item 能够同时提供精确的暗物质密度场和速度场，在保持大尺度结构与N体模拟一致的前提下，尽可能准确地刻画小尺度细节。
\end{itemize}

其中，星系巡天主要提取大尺度信息，而 CMB 透镜和弱引力透镜则需要精确的小尺度物质场。

\subsection{宇宙学快速模拟方法的研究现状与不足}

传统的快速模拟方法大多基于拉格朗日微扰论发展而来。例如，DESI DR1中使用的EZMock方法基于Zel'dovich近似（一阶拉格朗日微扰论），而DESI DR2计划采用的HOLI模拟则基于ALPT。微扰论方法计算简单、效率高，在宇宙早期过密度较小、引力演化近似线性的条件下是有效的；但在演化后期，引力非线性效应显著增强，该方法不再适用。因此，基于微扰论的快速模拟虽然能够快速生成大尺度上较为准确的结果，但在小尺度细节刻画方面仍有不足。经过进一步处理后，这类方法勉强能够满足当前星系巡天的需求，但难以满足下一代宇宙学探测对非线性小尺度精度的要求。

除微扰论外，也有研究尝试基于有效场论进行暗物质和星系分布的模拟，但有效场论存在计算速度较慢的问题，并不适合发展为快速模拟方法。

随着计算能力的提升和神经网络理论的发展，已有工作尝试使用UNet架构的神经网络来生成任意红移下的暗物质粒子的位移场和动量场，也有研究采用类似架构预测任意红移下的物质密度场。然而，这些工作通常采用较高的粒子分辨率，且大多未能同时给出粒子速度场的分布信息，速度场对于红移空间畸变、动理学SZ效应等探针的分析至关重要，而密度场和速度场的联合约束能力则是多探针宇宙学分析的重要优势。

\subsection{本文研究内容与贡献}

针对上述方法的不足，本文提出了一种基于生成对抗网络（GAN）的宇宙学快速模拟方法。该方法采用尺度分离的思想：在大尺度上，使用Zel'dovich近似快速生成位移场和速度场的线性大尺度结构；在小尺度上，训练两个独立的WGAN-GP神经网络分别学习位移场和速度场的残差，以补全非线性引力演化所导致的小尺度细节。该方法的主要贡献如下：
\begin{itemize}
  \item 同时生成暗物质密度场和速度场，密度场可服务于弱引力透镜等探针，速度场可服务于红移空间畸变等探针，两者结合可支持多探针联合分析；
  \item 采用残差学习策略，使神经网络仅需学习小尺度的非线性修正，降低了训练难度；
  \item 通过分块推理和旋转平均策略，使模型能够扩展到大规模模拟；
  \item 在单点统计、两点统计（功率谱）、三点统计（双谱）以及协方差矩阵等多个维度上，验证了该方法与N体模拟的一致性。
\end{itemize}

\section{文章结构}

本论文后续章节的组织结构如下：第2章详细阐述本文提出的快速模拟方法，包括整体工作流程、Zel'dovich近似生成大尺度结构、位移场与速度场残差网络的架构设计、损失函数以及训练与推理细节；第3章展示实验结果，从单点统计、功率谱、双谱和协方差矩阵等多个角度评估所提方法的性能；第4章总结全文工作，并讨论未来可能的研究方向。